%% =====================================================================
%%  Elementos pre-textuais
%%  Capa, folha de aprovacao, resumo, abstract, listas e sumario.
%%  Conforme NBR 14724 e NBR 6027.
%%
%%  Os campos entre colchetes sao [VALIDAR COM O AUTOR]: nao existiam no
%%  workspace e nao foram inventados. Ver MONOGRAFIA_PENDENCIAS.md.
%% =====================================================================

%% ---------------------------------------------------------------------
%%  CAPA
%% ---------------------------------------------------------------------
\begin{capa}
	\begin{center}

		\MakeUppercase{\TrabalhoInstituicao}

		\vspace{0.3cm}

		\MakeUppercase{\TrabalhoUnidade}

		\vspace{0.3cm}

		\MakeUppercase{\TrabalhoCurso}

		\vspace{2.5cm}

		\MakeUppercase{\TrabalhoAutor}

		\vspace{3.0cm}

		\textbf{\TrabalhoTitulo}

		\vfill

		\TrabalhoCidade

		\TrabalhoAno

	\end{center}
\end{capa}

%% ---------------------------------------------------------------------
%%  FOLHA DE APROVACAO
%% ---------------------------------------------------------------------
\begin{folhadeaprovacao}
	\begin{center}

		\MakeUppercase{\TrabalhoAutor}

		\vspace{2.0cm}

		\makebox[\textwidth]{\textbf{\TrabalhoTitulo}}

		\vspace{1.5cm}

		\begin{flushright}
			\begin{minipage}{8.5cm}
				\small\SingleSpacing
				\TrabalhoNatureza
				\vspace{0.4cm}

				Orientador: \TrabalhoOrientador
				\vspace{0.4cm}

				Coorientador: [VALIDAR COM O AUTOR --- ou declarar ausência]
			\end{minipage}
		\end{flushright}

		\vspace{1.5cm}

		Aprovado em: \rule{1cm}{0.4pt}/\rule{1cm}{0.4pt}/\rule{1.5cm}{0.4pt}

		\vspace{1.5cm}

		\rule{10cm}{0.4pt}\\
		\TrabalhoOrientador\\
		Presidente da banca --- \TrabalhoInstituicao

		\vspace{1.2cm}

		\rule{10cm}{0.4pt}\\
		\mbox{[VALIDAR COM O AUTOR --- Membro externo/interno 1]}\\
		\mbox{[VALIDAR COM O AUTOR --- titulação e instituição]}

		\vspace{1.2cm}

		\rule{10cm}{0.4pt}\\
		\mbox{[VALIDAR COM O AUTOR --- Membro externo/interno 2]}\\
		\mbox{[VALIDAR COM O AUTOR --- titulação e instituição]}

	\end{center}
\end{folhadeaprovacao}

%% ---------------------------------------------------------------------
%%  DEDICATORIA (opcional - remover se nao for usar)
%% ---------------------------------------------------------------------
%% \begin{dedicatoria}
%% 	\vspace*{\fill}
%% 	\begin{flushright}
%% 		\itshape [VALIDAR COM O AUTOR]
%% 	\end{flushright}
%% 	\vspace*{\fill}
%% \end{dedicatoria}

%% ---------------------------------------------------------------------
%%  AGRADECIMENTOS (opcional - remover se nao for usar)
%% ---------------------------------------------------------------------
%% \begin{agradecimentos}
%% 	[VALIDAR COM O AUTOR]
%% \end{agradecimentos}

%% ---------------------------------------------------------------------
%%  EPIGRAFE (opcional - remover se nao for usar)
%% ---------------------------------------------------------------------
%% \begin{epigrafe}
%% 	\vspace*{\fill}
%% 	\begin{flushright}
%% 		\itshape ``[VALIDAR COM O AUTOR --- citação com autoria]''
%% 	\end{flushright}
%% 	\vspace*{\fill}
%% \end{epigrafe}

%% ---------------------------------------------------------------------
%%  RESUMO - NBR 6028
%% ---------------------------------------------------------------------
\begin{resumo}
	Agentes de inteligência artificial generativa produzem código com alta
	velocidade, mas a ausência de um processo que force especificação prévia,
	critérios de aceitação verificáveis e registro de evidências compromete a
	auditabilidade e a manutenção do que foi produzido. Este trabalho define,
	aplica e avalia um processo de desenvolvimento orientado por especificações
	(\textit{spec-driven development} --- SDD) executado por agentes de
	inteligência artificial, organizado em quinze agentes especializados,
	quatro conjuntos de instruções permanentes e cinco habilidades de domínio,
	com fluxo adaptativo ao risco e seis portões de verificação. O processo foi
	aplicado à construção de um jogo digital completo entregue como artefato
	único de página web, especificado com quinze requisitos funcionais, nove
	requisitos não funcionais e doze critérios de aceitação. A verificação foi
	empírica, combinando inspeção estática do artefato com automação de
	navegador para teclado, apontador e leitura de pixels da área de desenho.
	Os doze critérios de aceitação foram satisfeitos, incluindo o bloqueio de
	inversão de trajetória, o sorteio de alvos fora do corpo do jogador e a
	persistência do recorde no armazenamento local. Mediu-se ainda o intervalo
	entre passos de simulação em duas faixas de pontuação: 133\,ms para
	pontuação inferior a quarenta pontos e 105\,ms para pontuação igual ou
	superior a oitenta pontos, sobre 184 movimentos amostrados, confirmando o
	aumento de velocidade previsto na especificação. O trabalho identifica como
	limitações a ausência de grupo de controle, a extensão reduzida da amostra
	temporal, a carência de critérios de aceitação para nove requisitos não
	funcionais e a não validação dos controles de toque em dispositivo físico.
	Conclui-se que o processo proposto torna cada entrega rastreável a um
	requisito e sustentada por evidência reproduzível, ao custo de maior
	disciplina documental.

	\vspace{\onelineskip}

	\noindent
	\textbf{Palavras-chave}: desenvolvimento orientado por especificações;
	agentes de inteligência artificial; engenharia de requisitos; verificação e
	validação; rastreabilidade; estudo de caso.
\end{resumo}

%% ---------------------------------------------------------------------
%%  ABSTRACT - NBR 6028
%% ---------------------------------------------------------------------
\begin{abstract}
	Generative artificial intelligence agents produce code at high speed, yet
	without a process that enforces prior specification, verifiable acceptance
	criteria and recorded evidence, the outcome is hard to audit and to
	maintain. This work defines, applies and evaluates a spec-driven
	development (SDD) process executed by artificial intelligence agents,
	organised into fifteen specialised agents, four permanent instruction sets
	and five domain skills, with a risk-adaptive flow and six verification
	gates. The process was applied to building a complete digital game
	delivered as a single web page artefact, specified with fifteen functional
	requirements, nine non-functional requirements and twelve acceptance
	criteria. Verification was empirical, combining static inspection of the
	artefact with browser automation for keyboard, pointer and pixel reading of
	the drawing surface. All twelve acceptance criteria passed, including
	trajectory reversal blocking, target spawning outside the player's body and
	high-score persistence in local storage. The interval between simulation
	steps was measured across two score bands: 133\,ms below forty points and
	105\,ms at or above eighty points, over 184 sampled movements, confirming
	the speed increase defined in the specification. The work identifies as
	limitations the absence of a control group, the small temporal sample, the
	lack of acceptance criteria for nine non-functional requirements and the
	unvalidated touch controls on physical devices. It concludes that the
	proposed process makes every deliverable traceable to a requirement and
	supported by reproducible evidence, at the cost of greater documentation
	discipline.

	\vspace{\onelineskip}

	\noindent
	\textbf{Keywords}: spec-driven development; artificial intelligence agents;
	requirements engineering; verification and validation; traceability; case
	study.
\end{abstract}

%% ---------------------------------------------------------------------
%%  LISTA DE ILUSTRACOES (figuras)
%% ---------------------------------------------------------------------
\pdfbookmark[0]{\listfigurename}{lof}
\listoffigures
\newpage

%% ---------------------------------------------------------------------
%%  LISTA DE QUADROS
%% ---------------------------------------------------------------------
\pdfbookmark[0]{Lista de quadros}{loq}
\listofquadro
\newpage

%% ---------------------------------------------------------------------
%%  LISTA DE TABELAS
%% ---------------------------------------------------------------------
\pdfbookmark[0]{\listtablename}{lot}
\listoftables
\newpage

%% ---------------------------------------------------------------------
%%  LISTA DE ABREVIATURAS E SIGLAS
%% ---------------------------------------------------------------------
\begin{siglas}
	\item[ABNT] Associação Brasileira de Normas Técnicas
	\item[ADR] \textit{Architecture Decision Record} (registro de decisão arquitetural)
	\item[API] \textit{Application Programming Interface}
	\item[AP] \textit{Access Point} (ponto de acesso sem fio)
	\item[CA] Critério de Aceitação
	\item[CSS] \textit{Cascading Style Sheets}
	\item[DPR] \textit{Device Pixel Ratio} (razão de pixels do dispositivo)
	\item[ECMA] \textit{European Computer Manufacturers Association}
	\item[FR] Requisito Funcional
	\item[FSM] \textit{Finite State Machine} (máquina de estados finitos)
	\item[HTML] \textit{HyperText Markup Language}
	\item[IA] Inteligência Artificial
	\item[ISO] \textit{International Organization for Standardization}
	\item[JSON] \textit{JavaScript Object Notation}
	\item[LLM] \textit{Large Language Model} (modelo de linguagem de grande porte)
	\item[NBR] Norma Brasileira
	\item[NFR] Requisito Não Funcional
	\item[RAM] \textit{Random Access Memory}
	\item[SDD] \textit{Spec-Driven Development}
	\item[SQuaRE] \textit{Software product Quality Requirements and Evaluation}
	\item[SRP] \textit{Single Responsibility Principle}
	\item[TCC] Trabalho de Conclusão de Curso
	\item[UI] \textit{User Interface} (interface do usuário)
	\item[V\&V] Verificação e Validação
	\item[W3C] \textit{World Wide Web Consortium}
	\item[WCAG] \textit{Web Content Accessibility Guidelines}
	\item[WHATWG] \textit{Web Hypertext Application Technology Working Group}
\end{siglas}

%% ---------------------------------------------------------------------
%%  LISTA DE SIMBOLOS
%% ---------------------------------------------------------------------
\begin{simbolos}
	\item[$A$] Conjunto de critérios de aceitação verificados
	\item[$C$] Custo agregado de verificação
	\item[$e$] Número de itens consumidos na partida
	\item[$\mathrm{E}{[}T{]}$] Intervalo médio entre passos de simulação
	\item[$m$] Número de amostras de movimento
	\item[$s$] Pontuação acumulada na partida
	\item[$t_{\max}$] Intervalo máximo entre passos de simulação
	\item[$t_{\min}$] Intervalo mínimo entre passos de simulação
	\item[$t_0$] Intervalo inicial entre passos de simulação
	\item[$\Delta$] Passo de redução do intervalo por item consumido
	\item[$\sigma$] Desvio padrão amostral
\end{simbolos}

%% ---------------------------------------------------------------------
%%  SUMARIO - NBR 6027
%% ---------------------------------------------------------------------
\pdfbookmark[0]{\contentsname}{toc}
\tableofcontents*
\newpage
